% Part: counterfactuals
% Chapter: introduction
% Section: counterfactuals

\documentclass[../../../include/open-logic-section]{subfiles}

\begin{document}

\olfileid{cnt}{int}{cnt}

\olsection{ప్రతివాస్తవ సోపాధికాలు}

ఆంగ్లంలోని ``ఒకవేళ \dots అయితే \dots''
నిర్మాణాల్లో చాలా సాధారణమైన, ముఖ్యమైన
ఒక రూపం \emph{to be} క్రియ యొక్క
భూతకాల సంభావనార్థక రూపాన్ని వాడుతుంది:
``ఒకవేళ \dots జరిగి ఉండి ఉంటే,
అప్పుడు \dots జరిగి ఉండేది.'' అటువంటి
సోపాధికం పూర్వపక్షం సాధారణంగా అసత్యం,
అంటే వాస్తవానికి విరుద్ధం కాబట్టి,
వాటిని \emph{ప్రతివాస్తవ సోపాధికాలు}
అంటారు. \emph{to be} క్రియ యొక్క
సంభావనార్థక రూపాన్ని వాడతాయి కాబట్టి
\emph{సంభావనార్థక సోపాధికాలు} అని
కూడా అంటారు. వీటిని ``ఒకవేళ
\dots జరుగుతుంటే, అప్పుడు \dots
జరుగుతుంది'' అనే రూపంలో ఉండే
\emph{సూచనాత్మక} సోపాధికాల నుంచి
వేరు చేస్తారు. ప్రతివాస్తవ,
సూచనాత్మక సోపాధికాల సత్య షరతులు
వేర్వేరు. ఆడమ్స్ ప్రసిద్ధ ఉదాహరణను
పరిశీలించండి:
\begin{quote}
  ఓస్వాల్డ్ కెన్నెడీని చంపకపోతే,
  మరెవరో చంపారు.
  
  ఓస్వాల్డ్ కెన్నెడీని చంపి
  ఉండకపోయి ఉంటే, మరెవరో
  చంపి ఉండేవారు.
\end{quote}
మొదటిది సూచనాత్మకం, రెండోది
ప్రతివాస్తవం. మొదటిది స్పష్టంగా
సత్యం: అధ్యక్షుడు జాన్ ఎఫ్. కెన్నెడీని
\emph{ఎవరో ఒకరు} చంపారని మనకు
తెలుసు. వారెన్ నివేదికకు విరుద్ధంగా,
ఆ వ్యక్తి లీ హార్వే ఓస్వాల్డ్ కాకపోతే,
కెన్నెడీని మరెవరో చంపారు. రెండోది
వేరే విషయం చెబుతుంది. ఓస్వాల్డ్
కెన్నెడీని చంపి ఉండకపోతే, అంటే
డల్లాస్‌లో కాల్పులు జరగకుండా ఉండి
ఉంటే లేదా అవి విఫలమై ఉంటే,
ఆ తరువాత చరిత్ర మరొక హత్యాయత్నం
విజయవంతమయ్యేలా సాగి ఉండేదని అది
వాదిస్తుంది. అది సత్యం కావాలంటే,
జేఎఫ్‌కే మరణాన్ని నిర్ధారించేందుకు
శక్తిమంతమైన వర్గాలు కుట్ర పన్నాయని
అనుకోవాలి (చాలామంది జేఎఫ్‌కే కుట్ర
సిద్ధాంతవాదులు నమ్మినట్లుగా).

\emph{సూచనాత్మక} సోపాధికాన్ని భౌతిక
సోపాధికం సరిగ్గా పట్టుకుంటుందా అనేది
ఇప్పటికీ చర్చనీయాంశమే. ప్రత్యేకించి,
ఆంగ్ల సూచనాత్మక సోపాధికాలకు భౌతిక
సోపాధికమే సత్య షరతులు ఇస్తుందనే
భావనతో సరిపడే విధంగా దాని
వైరుధ్యాభాసాలను ``వివరించగలమా''
అనేది ప్రశ్న. అయితే ప్రతివాస్తవ
సోపాధికాలను భౌతిక సోపాధికాలతో
సరిగ్గా సంకేతీకరించలేమన్నది
వివాదరహితం. ఎందుకంటే ప్రతివాస్తవాల
పూర్వపక్షాలు సాధారణంగా అసత్యాలే
అయినా, అసత్య పూర్వపక్షం గల ప్రతి
ప్రతివాస్తవ సోపాధికమూ సత్యం కాదు.
ఉదాహరణకు, వారెన్ నివేదికను మీరు
నమ్ముతూ, జేఎఫ్‌కే హత్యకు కుట్ర
లేదని భావిస్తే, ఆడమ్స్ ప్రతివాస్తవ
సోపాధికం దీనికి ఉదాహరణ.

కారణాత్మక వాదనలో ప్రతివాస్తవ
సోపాధికాలకు ముఖ్య పాత్ర ఉంది:
కారణ సంబంధాలను వ్యక్తీకరించడం
వాటి ప్రధాన వినియోగాల్లో ఒకటి.
ఉదాహరణకు, అగ్గిపుల్లను గీస్తే
అది మండుతుంది; దీనిని ``ఈ
అగ్గిపుల్లను గీసి ఉంటే అది
మండి ఉండేది'' అని వ్యక్తీకరించవచ్చు.
భౌతిక సోపాధికాలను, సాధారణంగా
సూచనాత్మక సోపాధికాలను, ఇలా
వాడలేం: ``అగ్గిపుల్ల గీయబడింది
$\lif$ అగ్గిపుల్ల మండుతుంది''
అనే వాక్యం అగ్గిపుల్లను అసలు
గీయకపోతే సత్యమే; గీసి ఉంటే
ఏమయ్యేదో దానితో సంబంధం లేదు.
ఇంకా విచిత్రంగా, ``అగ్గిపుల్ల
గీయబడింది $\lif$ అగ్గిపుల్ల పూల
గుత్తిగా మారుతుంది'' అనే వాక్యమూ
దాన్ని అసలు గీయకపోతే సత్యమే;
కానీ గీసి ఉంటే అది నిశ్చయంగా
పూల గుత్తిగా మారి ఉండేది కాదు.

ప్రతివాస్తవాలకు ఖచ్చితమైన సరైన
తర్కం ఏమిటన్నదానిపై ఇంకా చర్చ
జరుగుతోంది. స్టాల్నేకర్, లూయిస్
ప్రభావవంతమైన విశ్లేషణ ఇచ్చారు.
వారి ప్రకారం, ``ఒకవేళ~$S$ జరిగి
ఉండి ఉంటే~$T$ జరిగి ఉండేది''
అనే ప్రతివాస్తవ వాక్యం, వాస్తవ లోకం
ఉన్న తీరుకు అత్యంత సమీపంగా ఉండి,
అందులో~$S$ సత్యమయ్యే ప్రతివాస్తవ
పరిస్థితిలో (``సాధ్యలోకం'') $T$
సత్యమైనప్పుడే సత్యం. సాధ్యలోకాల
సత్తాశాస్త్రాన్ని ప్రస్తావిస్తుంది
కాబట్టి దీనిని ``సత్తాత్మక''
విశ్లేషణ అంటారు. ఇతర విశ్లేషణలు
షరతుపర సంభావ్యతలను లేదా విశ్వాస
సవరణ సిద్ధాంతాలను వాడతాయి.
ప్రతివాస్తవాల కోసం అనేక భిన్న
తర్కాలు ప్రతిపాదించబడ్డాయి.
ఒకే ఒక్క లూయిస్--స్టాల్నేకర్
ప్రతివాస్తవ తర్కమూ లేదు: ఇద్దరూ
సమీప సాధ్యలోకాలను ఆధారంగా
తీసుకొని ఒకే తరహా వివరణలు
ఇచ్చినప్పటికీ, వారి పరికల్పనలు
వేర్వేరు చెల్లుబాటు నిగమనాలకు
దారితీస్తాయి.

\end{document}
